% Figura 6.3: Frontera eficiente de Markowitz, cartera de mínima varianza, caso Coca-Cola vs Reebok y CML de Tobin
\begin{tikzpicture}
\begin{axis}[
    width=12.4cm,
    height=7.0cm,
    xlabel={Riesgo total: desviación típica $\sigma_p$ (\%)},
    ylabel={Rentabilidad esperada $\mathbb{E}[R_p]$ (\%)},
    xmin=0, xmax=65,
    ymin=0, ymax=24,
    xtick={0, 10, 20, 30, 40, 50, 60},
    xticklabels={{0\%}, {10\%}, {20\%}, {30\%}, {40\%}, {50\%}, {60\%}},
    ytick={0, 4, 8, 12, 16, 20, 24},
    yticklabels={{0\%}, {4\%}, {8\%}, {12\%}, {16\%}, {20\%}, {24\%}},
    grid=major,
    grid style={dotted, color=dfinmuted!45},
    axis line style={color=dfinslate, line width=0.9pt},
    tick label style={font=\scriptsize\sffamily, color=dfinslate},
    label style={font=\footnotesize\sffamily\bfseries, color=dfinnavy},
    legend style={at={(0.98,0.16)}, anchor=south east, font=\scriptsize\sffamily,
                  fill=white, draw=dfinslate!35, rounded corners=2pt}
]

% Curva exacta Coca-Cola (E=10%, sigma=31.5%) + Reebok (E=20%, sigma=58.5%), rho = 0.2
% Parámetro w_R en [0, 1]: E(w) = 10 + 10*w,
% sigma(w) = sqrt((1-w)^2*31.5^2 + w^2*58.5^2 + 2*(1-w)*w*0.2*31.5*58.5)
% Tramo completo de combinaciones (0 <= w <= 1):
\addplot[
    domain=0:1, samples=70, variable=\w,
    color=dfinmuted, line width=1.1pt, dashed
] (
    { sqrt((1-\w)^2*992.25 + \w^2*3422.25 + 2*(1-\w)*\w*368.55) },
    { 10 + 10*\w }
);
\addlegendentry{Combinaciones posibles ($\rho = +0{,}20$)}

% Rama superior eficiente desde cartera de mínima varianza global (w_min = (992.25-368.55)/(992.25+3422.25-737.1) = 623.7 / 3677.4 = 0.1696 => E=11.70%, sigma=29.07%) hasta w=1
\addplot[
    domain=0.1696:1.05, samples=60, variable=\w,
    color=dfinemerald, line width=1.8pt
] (
    { sqrt((1-\w)^2*992.25 + \w^2*3422.25 + 2*(1-\w)*\w*368.55) },
    { 10 + 10*\w }
);
\addlegendentry{Frontera eficiente con activos arriesgados}

% Punto Coca-Cola (w=0: sigma=31.5, E=10)
\addplot[only marks, mark=*, mark size=2.5pt, color=dfinnavy] coordinates {(31.5, 10.0)};
\node[font=\scriptsize\sffamily, color=dfinnavy, anchor=north west, fill=white, inner sep=1.5pt]
    at (axis cs:32.3, 9.7) {Coca-Cola ($31{,}5\%;\, 10\%$)};

% Punto Reebok (w=1: sigma=58.5, E=20)
\addplot[only marks, mark=*, mark size=2.5pt, color=dfinburg] coordinates {(58.5, 20.0)};
\node[font=\scriptsize\sffamily, color=dfinburg, anchor=north, fill=white, inner sep=1.5pt]
    at (axis cs:55.5, 19.3) {Reebok ($58{,}5\%;\, 20\%$)};

% Cartera 65% Coca-Cola / 35% Reebok (w=0.35: E=13.5%, sigma=31.72%)
\addplot[only marks, mark=square*, mark size=2.6pt, color=dfincopper] coordinates {(31.72, 13.5)};
\node[font=\scriptsize\sffamily\bfseries, color=dfincopper, anchor=west, fill=white, inner sep=1.5pt]
    at (axis cs:32.8, 13.5) {Cartera 65/35 ($31{,}7\%;\, 13{,}5\%$)};

% Cartera Mínima Varianza (29.07%, 11.70%)
\addplot[only marks, mark=diamond*, mark size=3pt, color=dfinemerald] coordinates {(29.07, 11.70)};
\node[font=\scriptsize\sffamily, color=dfinemerald, anchor=north east, fill=white, inner sep=1.5pt]
    at (axis cs:28.5, 11.45) {CMV};

% Línea del Mercado de Capitales (CML) desde rf = 3.5% tangente a w_T ≈ 0.56 (calculado numéricamente: rf=3.5 -> max Sharpe alrededor de w=0.53 -> sigma=36.93%, E=15.3%, pendiente (15.3-3.5)/36.93 = 0.3195)
% Recta desde (0, 3.5) hasta (62, 23.31):
\addplot[domain=0:62, samples=2, color=dfinnavy, line width=1.45pt]
    { 3.5 + 0.3195 * x };
\addlegendentry{CML (Línea del Mercado de Capitales)}

% Punto tangencia T y activo libre de riesgo
\addplot[only marks, mark=*, mark size=3pt, color=dfinnavy] coordinates {(0, 3.5) (36.93, 15.3)};
\node[font=\scriptsize\sffamily\bfseries, color=dfinnavy, anchor=south west, fill=white, inner sep=1.5pt]
    at (axis cs:1.2, 3.5) {$r_f = 3{,}5\%$};
\node[font=\scriptsize\sffamily\bfseries, color=dfinnavy, anchor=south east, fill=white, inner sep=1.5pt]
    at (axis cs:36.2, 16.0) {Cartera tangente $M$};

\end{axis}
\end{tikzpicture}
